% Ad Atticum 13.9 (ad-atticum-13-09) --- English translation
% Translated by Claude (Anthropic) from the Latin (Perseus).
% Style guide: STYLE.md.

\ciceroLetterOpener{CICERO ATTICO salutem}

\ciceroSection{1} You had only just left yesterday when Trebatius arrived, and a little later Curtius --- the latter to pay his respects, but he stayed when invited. Trebatius I am keeping with me. This morning, Dolabella. Long conversation, far into the day. I cannot describe anything more attentive \textit{[Greek: ektenesteron]}, anything more affectionate \textit{[Greek: philostorgoteron]}. It came round, however, to Quintus. Much that is unspeakable \textit{[Greek: aphata]}, not to be recounted \textit{[Greek: adiegeta]}, but one thing of such a kind that, unless the army knew it, I should not dare even to dictate it to Tiro, let alone write it myself --- but let that suffice. Just at the right moment \textit{[Greek: eukairos]}, while I had Dolabella with me, Torquatus came to me, and most courteously Dolabella set out for him the terms in which I had taken up his case. As it happened I had taken it up most carefully; and the care looked welcome to Torquatus.

\ciceroSection{2} I wait to hear from you if there is anything about Brutus. Though Nicias was thinking the thing settled, with the divorce only failing to find approval --- about which I am all the more anxious, on the same grounds you are. For if there is anything in the way of a quarrel, this step can heal it. I must go to Arpinum --- those small estates of mine need to be put in order by me, and I am afraid there will be no chance of getting away once Caesar arrives. On his coming, Dolabella holds the same view that you were inferring from Messalla's letter. When I have got there and made out what the business is, I will write to you about the days for my return.
